\paragraph{Content class of the order-\(\LUV^2\) comparison: the gaps at general content, the ledger at \(N_s=6\), \(n_D=2\) and sectors with a scale (conditional; not closure).}
Keep the setting of the paragraphs on the normalisation of the tensor kinetic term, on fermionic content and on the order \(\Lambda^2\): four matrix coordinates, the Poisson limit taken before the one common proper-time limit, the Euclidean section, \(\alpha=1\), one loop, the \(\LUV^2\) term, \(N_s\) scalar components, the connection with its ghosts and \(n_D\) Dirac fields with the Yukawa couplings dropped, and the two modes with \(\delta\rho=0\) about a constant \(\bar\Theta\). No state of the matrix model enters: the replicas are the gluing prescriptions of the preceding paragraphs for the Poisson-limit operators on the cover of the matrix chart, and no limit of finite matrices is taken.

\emph{Gaps at general content.} The sector densities \(\tfrac n6\rho R[g_\rho]+\kappa\,\rho^2X\), with \((n,\kappa)=(1,0)\) for a scalar component and a ghost, \((4,1)\) for the connection and \((4,\tfrac12)\) for a Dirac field, and the eigenvalue \(-\tfrac14\) of \(\rho R[g_\rho]\) relative to \(\rho^2X\) on the modes with \(\delta\rho=0\), are those of the preceding paragraphs (checked from \(\mathcal O_s\), \(\mathcal O_v\), \(\mathcal O_f\) at four random two-jets of a closed \(\omega\) and on four random modes, exact arithmetic). With \(c_R=N_{\mathrm{tot}}/6\), \(c_T=c_R-4\kappa_X\), \(c_W(r)=c_R+\kappa_Xr\) and the area coefficients \(c_A^{\mathrm{per}}=c_R\) of the unrotated periodic gluing and \(c_A=c_R+n_D\) of the fixed-state replica, for every \(N_s\) and \(n_D\)
\[
  c_A-c_T=4\kappa_X+n_D,\qquad c_A^{\mathrm{per}}-c_T=4\kappa_X,\qquad (c_A-1)-c_W(-1)=\tfrac12n_D,\qquad (c_A-1)-c_T=3\kappa_X+\tfrac12n_D .
\]
The first two are the gaps \(4-n_D\) and \(4-2n_D\) of the paragraph on the order \(\Lambda^2\). The identification of the fixed-state gap with \(4\kappa_X\), the relation \(c_A-1=c_W(-1)\) and the remainder \(3\kappa_X\) of the paragraphs on covariant variables and on the matrix gauge realisation hold at \(n_D=0\), as stated there, and at no other \(n_D\). The difference \(c_T-c_W(0)=-4\kappa_X\) does not depend on \(n_D\), but for \(n_D>0\) the fixed-state coefficient exceeds \(c_W(0)\) by the Dirac excess \(n_D\), which lies on the area side only; the term that the fixed-state replica needs is then \(-(4\kappa_X+n_D)\), and the coefficient \(-4\kappa_X\) required at \(n_D=0\) in the paragraphs on the matrix gauge realisation and on the nearest published statements is its value at \(n_D=0\).

\emph{At \(N_s=6\), \(n_D=2\).} Here \(c_R=\kappa_X=0\) and the weighted \(a_1\) density vanishes pointwise (paragraph on the next order), so \(c_T=c_T^{\mathrm{ext}}=c_W(r)=0\) for every \(r\): the family \(L_r\) of the paragraph on covariant variables reduces to \(L_0\), and the presentation dependence of the Noether coefficient and the coefficient \(-4\kappa_X\) vanish. The periodic and smooth replicas give \(0=0\), and the uniform contact \(-N_{\mathrm{tot}}/4\) is zero. The one nonvanishing entry is the fixed-state coefficient \(c_A=n_D=2\), the Dirac excess alone, so the term that the fixed-state replica needs is \(-2\) and is not a term in \(\rho^2X\). The Dirac holonomy contact vanishes at first order, and the vector holonomy contact of modulus \(1\) changes \(c_A\) to \(1\) on \(n<1\), the value \(-1\) being that of the edge-mode term of~\cite{DonnellyWall2016} for a Maxwell field if the state carries that sector (an input, not a prescription of this class), and to \(3\) on \(n>1\); no contact of the paragraph on the order \(\Lambda^2\) removes it. At logarithmic order, about a constant self-dual or anti-self-dual \(\bar\Theta\) on the \(n<1\) branch, the gaps of the coefficient of \(W\) of the paragraph on the content dependence of the logarithmic comparison are \(0\), \(2\), \(0\) at \(n_D=2\), with \(\kappa_C^T=0\), and the logarithmic coefficient \(\propto\int\rho^2\mathcal S_4\) of the paragraph on the next order, quartic in \(b_\sigma\), vanishes at quadratic order about a constant \(\bar\Theta\), given the sampled-jet identity of that paragraph. In the units \(1/G=c\,\LUV^2/2\pi\) of the paragraph on the order \(\Lambda^2\), the requirement \(G_A=G_T\) of the paragraph on the normalisation is, at this content and order, the single condition \(c_A=0\) with \(1/G_T=0\): for the fixed-state replica, which is the trace of \(\rho^n\) only under the hypotheses of the paragraph on fermion replica gluing (a quasi-free state of the Poisson-limit fermions with rotational modular flow, the less singular extension and a unique continuation from integer \(n\)), \(1/G_A=\LUV^2/\pi\) while \(1/G_T=0\), and for the periodic and smooth replicas both vanish. In this class the comparison at this content is an equality of vanishing coefficients or the inequality \(2\neq0\), and in neither case does it determine a Newton constant.

\emph{Sectors with a scale.} Let each sector carry a constant mass, \(\mathcal O_i\to\mathcal O_i+m_i^2\), the formal shift of the paragraph on the next order, with the ghosts at the mass of the connection. The heat trace of the sector is multiplied by \(e^{-sm_i^2}\) on the chart and on the cover alike, so both sides are weighted by \(f(m^2)=\int_{\LUV^{-2}}^\infty\mathrm ds\,s^{-2}e^{-sm^2}=\LUV^2-m^2\ln(\LUV^2/m^2)-(1-\gamma_{\mathrm E})m^2+O(m^4/\LUV^2)\), with \(\gamma_{\mathrm E}\) the Euler constant; on the tensor side the term in \(m^2\ln\LUV\) is the curvature part of the change of \(a_2\) stated there. With \(\mathsf S_n=\sum_iw_in_if(m_i^2)\) and \(\mathsf S_\kappa=\sum_iw_i\kappa_if(m_i^2)\), in units in which the massless value of \(f\) is \(1\),
\[
  c_T=\tfrac16\mathsf S_n-4\mathsf S_\kappa,\qquad c_A^{\mathrm{per}}=\tfrac16\mathsf S_n,\qquad c_A=\tfrac16\mathsf S_n+\sum_{\text{Dirac}}f(m_D^2),\qquad c_A^{\mathrm{per}}-c_T=4\mathsf S_\kappa .
\]
At \(N_s=6\), \(n_D=2\) the \(\LUV^2\) parts of \(\mathsf S_n\) and \(\mathsf S_\kappa\) cancel and both sums are of order \(m^2\ln\LUV\). The tensor coefficient is nonzero only if \(\mathsf S_n\neq24\,\mathsf S_\kappa\), and for the periodic gluing it equals the area coefficient if and only if \(\mathsf S_\kappa=0\), a condition on the connection and Dirac spectra alone; the fixed-state replica keeps in addition the Dirac term, of order \(\LUV^2\). For one bosonic mass and one fermionic mass \(c_T=-\tfrac83(f_b-f_f)\) and \(c_A^{\mathrm{per}}=\tfrac43(f_b-f_f)\), of opposite sign. For this formal shift the comparison at this content is thus not vacuous only for a split spectrum, where it reduces to the two sums \(\mathsf S_n\) and \(\mathsf S_\kappa\). The product backgrounds \(M\times K\) of the paragraph on product backgrounds in Section~\ref{sec:status-F} are not of this form: there the scalar operator of \(K\) is common to all species, the term that would carry \(\mathsf S_\kappa\) has no two-derivative part, and the comparison is \(c_T=\tfrac16W_R-4W_X\) against \(c_A^{\mathrm{per}}=\tfrac16W_R\), with \(W_X\) from the mixed term, which that paragraph relates to a term of Section~6.2 of~\cite{Steinacker2023}, so that \(W_R\) and \(W_X\) take the places of \(\mathsf S_n\) and \(\mathsf S_\kappa\) and the periodic gap \(4W_X\) is nonzero for a nonzero flux of \(K\); the fixed-state replica there keeps the value \(c_A=2\) of this content for each matrix mode of \(K\) at order \(\LUV^2\), with cutoff-dependent terms whose finite remainders are not computed there. Those statements are for the Poisson limit on \(M\), one loop, a fuzzy space \(K\) that is not constructed and not required to solve the equations of motion, product backgrounds without mixed flux, the orientation-even part of the mixed term and the replicas of this section, and they fix no sign of \(W_R\).

These statements are for the Poisson limit taken first, four matrix coordinates, the Euclidean section, one loop, \(\alpha=1\), one common proper-time limit, the contents stated with the Yukawa couplings dropped, the replicas and contacts of the preceding paragraphs with their hypotheses, the two modes with \(\delta\rho=0\), and formal constant masses without a gauge-invariant mass term. They supply no state, modular Hamiltonian or gluing of the matrix model, do not derive the content, treat no finite terms, Moyal corrections or Kaluza--Klein tower, add nothing to the statements quoted from Section~\ref{sec:status-F}, and give no value or sign of \(G_{\mathrm{eff}}\). Problem~B stays open and PROGRAMME-CONDITIONAL.

